\chapter{Complexity — NARM}
%%% generated by the blueprint pipeline from TCSlib.Complexity.SpaceComplexity.Machines.NARM
\section{Overview}

Declarations of the Lean module \texttt{TCSlib.Complexity.SpaceComplexity.Machines.NARM}.

\section{Declarations}

\begin{definition}[Nondeterministic register instructions]
\lean{Complexity.LogProg.NIns}
\label{Complexity.LogProg.NIns}
\leanok
\uses{Complexity.LogProg.Ins}
For $m$ registers, $d$ deciders, and a set $\Lambda$ of control states, a
nondeterministic register instruction is either a deterministic register instruction or
a choice instruction with two successor states $\ell_1,\ell_0\in\Lambda$, selected
according to whether the choice bit is $1$ or $0$.
\end{definition}

\begin{definition}[Nondeterministic register program]
\lean{Complexity.LogProg.NARM}
\label{Complexity.LogProg.NARM}
\leanok
\uses{Complexity.LogProg.NIns}
For $m$ registers, $d$ deciders, and a set $\Lambda$ of control states, a
nondeterministic register program assigns a nondeterministic register instruction to
each control state in $\Lambda$.
\end{definition}

\begin{definition}[Successor targets of an instruction]
\lean{Complexity.LogProg.NIns.targets}
\label{Complexity.LogProg.NIns.targets}
\leanok
\uses{Complexity.LogProg.NIns}
For a nondeterministic register instruction, define its target pair to be undefined if
it is deterministic, and to be $(\ell_1,\ell_0)$ if it is a choice instruction with
successors $\ell_1$ and $\ell_0$.
\end{definition}

\begin{definition}[Deterministic part of a nondeterministic program]
\lean{Complexity.LogProg.NARM.detPart}
\label{Complexity.LogProg.NARM.detPart}
\leanok
\uses{Complexity.LogProg.ARM, Complexity.LogProg.NARM}
Given a nondeterministic register program and a Boolean value $\beta$, its deterministic
part assigns each control state its deterministic instruction when the original
instruction is deterministic; when it is a choice instruction, it instead assigns the
instruction that returns $\beta$ and halts.
\end{definition}

\begin{lemma}[Deterministic Part at a Deterministic Instruction]
\lean{Complexity.LogProg.NARM.detPart_det}
\label{Complexity.LogProg.NARM.detPart_det}
\leanok
\uses{Complexity.LogProg.Ins, Complexity.LogProg.NARM, Complexity.LogProg.NARM.detPart}
\difficulty{1}
Let $A$ be a nondeterministic register machine, let $l$ be a state, and suppose the
instruction of $A$ at $l$ is the deterministic instruction $i$. For either Boolean value
$\beta$, the deterministic part of $A$ at $l$, formed using $\beta$, is $i$.
\end{lemma}

\begin{lemma}[Deterministic Part at a Choice Instruction]
\lean{Complexity.LogProg.NARM.detPart_choose}
\label{Complexity.LogProg.NARM.detPart_choose}
\leanok
\uses{Complexity.LogProg.NARM, Complexity.LogProg.NARM.detPart}
\difficulty{1}
Let $A$ be a nondeterministic register machine, let $l$ be a state, and suppose the
instruction of $A$ at $l$ is a choice between states $l_1$ and $l_0$. For either Boolean
value $\beta$, the deterministic part of $A$ at $l$, formed using $\beta$, is the
instruction that returns $\beta$ and halts.
\end{lemma}

\begin{definition}[One nondeterministic computation step]
\lean{Complexity.LogProg.nastep}
\label{Complexity.LogProg.nastep}
\leanok
\uses{Complexity.LogProg.AConf, Complexity.LogProg.NARM, Complexity.LogProg.astep}
Given a nondeterministic register program, an oracle, an input bit string, and a choice
bit $b$, one step leaves a configuration with no active control state unchanged. At an
active state it executes the associated deterministic instruction, or, for a choice
instruction with successors $\ell_1,\ell_0$, moves to $\ell_1$ if $b=1$ and to $\ell_0$
if $b=0$, leaving the registers and result unchanged.
\end{definition}

\begin{definition}[Run on a sequence of choice bits]
\lean{Complexity.LogProg.narun}
\label{Complexity.LogProg.narun}
\leanok
\uses{Complexity.LogProg.AConf, Complexity.LogProg.NARM, Complexity.LogProg.nastep}
Given a nondeterministic register program, an oracle, and an input bit string, its run
on a finite sequence of choice bits is obtained by processing those bits in order,
applying one nondeterministic computation step for each bit; the empty sequence leaves
the initial configuration unchanged.
\end{definition}

\begin{definition}[Nondeterministic precondition predicate]
\lean{Complexity.LogProg.NPreS}
\label{Complexity.LogProg.NPreS}
\leanok
\uses{Complexity.LogProg.AConf, Complexity.LogProg.NARM, Complexity.LogProg.NARM.detPart, Complexity.LogProg.PreS}
Given a nondeterministic register program $A$, an input bit string $x$, and a
register-machine configuration $a$, the predicate $\mathrm{NPreS}$ holds exactly when
the deterministic program obtained from $A$ by replacing each choice with the answer $0$
satisfies $\mathrm{PreS}$ on input $x$ at configuration $a$.
\end{definition}

\begin{definition}[Deterministic precondition of a nondeterministic program]
\lean{Complexity.LogProg.NPre}
\label{Complexity.LogProg.NPre}
\leanok
\uses{Complexity.LogProg.AConf, Complexity.LogProg.NARM, Complexity.LogProg.NARM.detPart, Complexity.LogProg.Pre, Turing.MultiTapeTM}
Given a nondeterministic register program $A$, an oracle, a deterministic multitape
Turing machine $D$, a state assignment for the oracle deciders, a bound $B$, an input
string $x$, and a register configuration $a$, the predicate holds exactly when the
deterministic precondition holds for the program obtained by resolving every
nondeterministic choice as $0$, with the same oracle, machine, state assignment, bound,
input, and configuration.
\end{definition}

\begin{lemma}[Deterministic Instruction Step]
\lean{Complexity.LogProg.nastep_det}
\label{Complexity.LogProg.nastep_det}
\leanok
\uses{Complexity.LogProg.Ins, Complexity.LogProg.NARM.detPart, Complexity.LogProg.NARM.detPart_det, Complexity.LogProg.astep, Complexity.LogProg.nastep}
\difficulty{3}
Let $A$ be a nondeterministic register program, let an oracle assign a Boolean answer to
each decider and bit string, and fix an input word, a program state $l$, and a
deterministic instruction $i$. If $A$ assigns the deterministic instruction $i$ to $l$,
then for either Boolean choice parameter $eta$ and any choice bit $b$, register
valuation, and optional result, one step of the nondeterministic program from state $l$
agrees with one step of the deterministic program obtained from $A$ by replacing each
choice instruction with the answer $eta$.
\end{lemma}

\begin{lemma}[Step at a choice instruction]
\lean{Complexity.LogProg.nastep_choose}
\label{Complexity.LogProg.nastep_choose}
\leanok
\uses{Complexity.LogProg.nastep}
\difficulty{3}
Let $A$ be a nondeterministic register machine, let $\text{oracle}$ be an oracle, let
$x$ be an input bit string, and suppose the instruction at control state $l$ chooses
between states $l_1$ and $l_0$. For any Boolean choice $b$, register assignment $v$, and
optional Boolean result $r$, one step from the configuration with control state $l$,
register assignment $v$, and result $r$ moves to $l_1$ if $b$ is true and to $l_0$ if
$b$ is false, leaving the registers and result unchanged.
\end{lemma}

\begin{lemma}[Halted Configuration Is Fixed]
\lean{Complexity.LogProg.nastep_none}
\label{Complexity.LogProg.nastep_none}
\leanok
\uses{Complexity.LogProg.nastep}
\difficulty{1}
Let $m,d$ be natural numbers, let $\Lambda$ be a type, let $A$ be a nondeterministic
register program with states in $\Lambda$, let $O$ assign a Boolean answer to each
of the $d$ deciders on each bit string, and let $x$ be a bit string. For every choice
bit $b$, register valuation $v : \\mathrm{Fin}(m) \\to \\mathbb{N}$, and optional
Boolean result $r$, one step of the program on input $x$ with oracle $O$ leaves the
halted configuration unchanged: $(\mathrm{none},v,r) \\mapsto (\mathrm{none},v,r)$.
\end{lemma}

\begin{lemma}[Empty choice sequence leaves the configuration unchanged]
\lean{Complexity.LogProg.narun_nil}
\label{Complexity.LogProg.narun_nil}
\leanok
\uses{Complexity.LogProg.AConf, Complexity.LogProg.narun}
\difficulty{1}
For any nondeterministic register machine $A$, oracle, input bit string $x$, and
configuration $a$, running $A$ on the empty choice sequence with oracle and input $x$
returns $a$.
\end{lemma}

\begin{lemma}[One-step unfolding of a run]
\lean{Complexity.LogProg.narun_cons}
\label{Complexity.LogProg.narun_cons}
\leanok
\uses{Complexity.LogProg.AConf, Complexity.LogProg.narun, Complexity.LogProg.nastep}
\difficulty{1}
Let $A$ be a nondeterministic register machine, let $\text{oracle}$ be an oracle, let
$x$ be an input bit string, let $b$ be a Boolean choice, let $\sigma$ be a finite choice
string, and let $a$ be an initial configuration. Running on the string formed by placing
$b$ before $\sigma$ is the same as first taking the step determined by $b$ and then
running on $\sigma$.
\end{lemma}

\begin{lemma}[Concatenation law for nondeterministic runs]
\lean{Complexity.LogProg.narun_append}
\label{Complexity.LogProg.narun_append}
\leanok
\uses{Complexity.LogProg.AConf, Complexity.LogProg.narun}
\difficulty{4}
Let $A$ be a nondeterministic register machine, let $	ext{oracle}$ be an oracle, let $x$
be an input bit string, and let $a$ be a machine configuration. For any two finite
choice strings $\sigma$ and $\tau$, running on their concatenation gives the same
configuration as first running on $\sigma$ and then running on $\tau$.
\end{lemma}

\begin{lemma}[Runs remain halted after halting]
\lean{Complexity.LogProg.narun_halted}
\label{Complexity.LogProg.narun_halted}
\leanok
\uses{Complexity.LogProg.narun}
\difficulty{4}
Let $A$ be a nondeterministic register machine, let $\text{oracle}$ be an oracle, let
$x$ be an input bit string, let $v$ assign a natural-number value to each of the $m$
registers, and let $r$ be an optional Boolean result. Running from the halted
configuration with register values $v$ and result $r$ leaves that configuration
unchanged for every finite choice string.
\end{lemma}

\begin{lemma}[Appending choices after halting]
\lean{Complexity.LogProg.narun_append_of_halted}
\label{Complexity.LogProg.narun_append_of_halted}
\leanok
\uses{Complexity.LogProg.AConf, Complexity.LogProg.narun, Complexity.LogProg.narun_append, Complexity.LogProg.narun_halted}
\difficulty{4}
Let $A$ be a nondeterministic register machine, let $	ext{oracle}$ be an oracle, let $x$
be an input bit string, and let $a$ be an initial configuration. If the run on a finite
choice string $\sigma$ has halted, then for every finite choice string $\tau$, running
on the concatenation $\sigma\tau$ yields exactly the same configuration as running on
$\sigma$.
\end{lemma}

\begin{lemma}[Choice-independent precondition]
\lean{Complexity.LogProg.pre_detPart}
\label{Complexity.LogProg.pre_detPart}
\leanok
\uses{Complexity.LogProg.AConf, Complexity.LogProg.NARM.detPart, Complexity.LogProg.NARM.detPart_choose, Complexity.LogProg.NARM.detPart_det, Complexity.LogProg.NPre, Complexity.LogProg.Pre, Turing.MultiTapeTM}
\difficulty{4}
For every nondeterministic program, oracle, decider, initial-state assignment, space
bound, input, Boolean replacement for choice instructions, and configuration, the
precondition of the deterministic program formed by replacing each choice with that
Boolean agrees with the nondeterministic precondition.
\end{lemma}

\begin{lemma}[Agreement with the deterministic simulation]
\lean{Complexity.LogProg.narun_eq_arun}
\label{Complexity.LogProg.narun_eq_arun}
\leanok
\uses{Complexity.LogProg.AConf, Complexity.LogProg.NARM.detPart, Complexity.LogProg.NIns.targets, Complexity.LogProg.arun, Complexity.LogProg.arun_succ, Complexity.LogProg.astep, Complexity.LogProg.narun, Complexity.LogProg.narun_cons, Complexity.LogProg.nastep, Complexity.LogProg.nastep_det}
\difficulty{6}
Let $A$ be a nondeterministic register machine, let $\text{oracle}$ be an oracle, let
$x$ be an input bit string, let $\beta$ be a Boolean value, and let $a$ be an initial
configuration. Suppose that at every time $t<T$, whenever the deterministic simulation
that replaces each choice instruction by the answer $\beta$ is in a control state $l$,
the instruction of $A$ at $l$ is not a choice instruction. Then for every finite choice
string $\sigma$ of length at most $T$, the nondeterministic run on $\sigma$ agrees with
that deterministic simulation after $|\sigma|$ steps.
\end{lemma}

\begin{definition}[Halting on every choice string]
\lean{Complexity.LogProg.NHalts}
\label{Complexity.LogProg.NHalts}
\leanok
\uses{Complexity.LogProg.AConf, Complexity.LogProg.NARM, Complexity.LogProg.narun}
Let $A$ be a nondeterministic register program, let an oracle and a Boolean input string
$x$ be fixed, and let $G$ be a predicate on configurations. For a natural number $N$,
every choice string of length $N$ makes the run from the given configuration halt after
those choices; moreover, at every time $t<N$ for which the run using the first $t$
choices has not halted, its configuration satisfies $G$.
\end{definition}

\begin{definition}[Acceptance with a fixed number of choices]
\lean{Complexity.LogProg.NAns}
\label{Complexity.LogProg.NAns}
\leanok
\uses{Complexity.LogProg.AConf, Complexity.LogProg.NARM, Complexity.LogProg.narun}
Let $A$ be a nondeterministic register program, let an oracle and a Boolean input string
$x$ be fixed, and let $a$ be an initial configuration. For a natural number $N$ and a
Boolean value $b$, there is a choice string of length $N$ whose run from $a$ produces
answer $b$.
\end{definition}

\begin{lemma}[Splitting a list prefix]
\lean{Complexity.LogProg.take_split}
\label{Complexity.LogProg.take_split}
\leanok
\difficulty{3}
Let $\sigma$ be a finite Boolean list and let $N,s$ be natural numbers with $N\le s$.
Then the prefix of $\sigma$ of length $s$ is the concatenation of its prefix of length
$N$ and the prefix of length $s-N$ of the list obtained by removing its first $N$
entries.
\end{lemma}

\begin{lemma}[Runs from an already halted configuration]
\lean{Complexity.LogProg.NHalts.halted}
\label{Complexity.LogProg.NHalts.halted}
\leanok
\uses{Complexity.LogProg.NHalts}
\difficulty{1}
For any nondeterministic register program, oracle, input, predicate $G$, register
valuation $v$, result value $r$, and natural number $N$, every choice word of length $N$
from the halted configuration $(\mathrm{none},v,r)$ leads to a halted configuration, and
every earlier configuration satisfies $G$ whenever it is nonhalted.
\end{lemma}

\begin{lemma}[Monotonicity of the halting bound]
\lean{Complexity.LogProg.NHalts.mono}
\label{Complexity.LogProg.NHalts.mono}
\leanok
\uses{Complexity.LogProg.AConf, Complexity.LogProg.NHalts}
\difficulty{1}
Let $A$ be a nondeterministic register program, with fixed oracle, input, and predicate
$G$. If every choice word of length $N$ from a configuration $a$ leads to a halted
configuration and all earlier nonhalted configurations satisfy $G$, then the same holds
for every choice word of length $N'$ whenever $N\le N'$.
\end{lemma}

\begin{lemma}[Monotonicity in choice-word length]
\lean{Complexity.LogProg.NHalts.mono_len}
\label{Complexity.LogProg.NHalts.mono_len}
\leanok
\uses{Complexity.LogProg.AConf, Complexity.LogProg.NHalts, Complexity.LogProg.narun, Complexity.LogProg.narun_append_of_halted, Complexity.LogProg.take_split}
\difficulty{5}
Let $A$ be a nondeterministic register program, with fixed oracle, input, and predicate
$G$. If every choice word of length $N$ from a configuration $a$ leads to a halted
configuration and all earlier nonhalted configurations satisfy $G$, then the same holds
for every choice word of length $N'$ whenever $N\le N'$.
\end{lemma}

\begin{lemma}[Halting after a deterministic instruction]
\lean{Complexity.LogProg.NHalts.det}
\label{Complexity.LogProg.NHalts.det}
\leanok
\uses{Complexity.LogProg.Ins, Complexity.LogProg.NARM.detPart, Complexity.LogProg.NHalts, Complexity.LogProg.astep, Complexity.LogProg.narun_cons, Complexity.LogProg.nastep_det}
\difficulty{4}
Let $A$ be a nondeterministic register program, with fixed oracle, input, and
configuration predicate $G$. Suppose the instruction at state $l$ is the deterministic
instruction $i$, and let $\beta$ be either Boolean value. If $(l,v,r)$ satisfies $G$,
and every choice word of length $N$ from the configuration obtained by applying $i$ at
$(l,v,r)$ leads to a halted configuration while all its earlier nonhalted configurations
satisfy $G$, then every choice word of length $N+1$ from $(l,v,r)$ has those same
halting and intermediate-configuration properties.
\end{lemma}

\begin{lemma}[Halting after a nondeterministic choice]
\lean{Complexity.LogProg.NHalts.choose}
\label{Complexity.LogProg.NHalts.choose}
\leanok
\uses{Complexity.LogProg.NHalts, Complexity.LogProg.narun_cons, Complexity.LogProg.nastep_choose}
\difficulty{5}
Let $A$ be a nondeterministic register program, with fixed oracle, input, and
configuration predicate $G$. Suppose its instruction at state $l$ chooses between states
$l_1$ and $l_0$. If the current configuration $(l,v,r)$ satisfies $G$, and every choice
word of length $N$ from each of $(l_1,v,r)$ and $(l_0,v,r)$ leads to a halted
configuration while all its earlier nonhalted configurations satisfy $G$, then every
choice word of length $N+1$ from $(l,v,r)$ has those same halting and
intermediate-configuration properties.
\end{lemma}

\begin{lemma}[Halting Along a Deterministic Run]
\lean{Complexity.LogProg.NHalts.run}
\label{Complexity.LogProg.NHalts.run}
\leanok
\uses{Complexity.LogProg.AConf, Complexity.LogProg.NARM.detPart, Complexity.LogProg.NHalts, Complexity.LogProg.NIns.targets, Complexity.LogProg.arun, Complexity.LogProg.narun, Complexity.LogProg.narun_append, Complexity.LogProg.narun_eq_arun, Complexity.LogProg.take_split}
\difficulty{5}
Let $A$ be a nondeterministic register machine, with oracle, input bit string, and
configuration guard $G$ fixed. Suppose that $T$ steps of the machine obtained by
replacing every choice instruction by the fixed return value $eta$ take configuration
$a$ to $a'$. At every time $t<T$, suppose the configuration reached from $a$ satisfies
$G$, and whenever it is live its current instruction in $A$ is deterministic. If the
halting condition holds for $N$ steps from $a'$, then it holds for $T+N$ steps from $a$.
\end{lemma}

\begin{lemma}[Initial guard condition]
\lean{Complexity.LogProg.NHalts.guard}
\label{Complexity.LogProg.NHalts.guard}
\leanok
\uses{Complexity.LogProg.NHalts}
\difficulty{3}
For any nondeterministic register program, oracle, input, and predicate $G$, if every
choice word of length $N+1$ from the configuration $(l,v,r)$ leads to a halted
configuration and all earlier nonhalted configurations satisfy $G$, then $G$ holds of
$(l,v,r)$.
\end{lemma}

\begin{lemma}[Halting After a Deterministic Step]
\lean{Complexity.LogProg.NHalts.of_det}
\label{Complexity.LogProg.NHalts.of_det}
\leanok
\uses{Complexity.LogProg.Ins, Complexity.LogProg.NARM.detPart, Complexity.LogProg.NHalts, Complexity.LogProg.astep, Complexity.LogProg.narun_cons, Complexity.LogProg.nastep_det}
\difficulty{4}
Let $A$ be a nondeterministic register machine, with oracle, input bit string, and
configuration guard $G$ fixed. Suppose the instruction at control state $l$ is the
deterministic instruction $i$, and the halting condition holds for $N+1$ steps from the
configuration at $l$. Then it holds for $N$ steps from the configuration obtained by
applying one step of the machine in which every choice instruction is replaced by the
fixed return value $eta$.
\end{lemma}

\begin{lemma}[Halting After a Choice]
\lean{Complexity.LogProg.NHalts.of_choose}
\label{Complexity.LogProg.NHalts.of_choose}
\leanok
\uses{Complexity.LogProg.NHalts, Complexity.LogProg.narun_cons, Complexity.LogProg.nastep_choose}
\difficulty{4}
Let $A$ be a nondeterministic register machine, let an oracle and input bit string be
fixed, and let $G$ be a predicate on configurations. If the instruction at control state
$l$ chooses between $l_1$ and $l_0$, then for either choice bit $b$, halting within the
stated bound from the configuration at $l$ implies halting within one fewer step from
the selected successor configuration. Here the halting condition requires every choice
sequence of the specified length to end in a halted configuration and every earlier live
configuration to satisfy $G$.
\end{lemma}

\begin{lemma}[Answer from a Halted Configuration]
\lean{Complexity.LogProg.nans_halted}
\label{Complexity.LogProg.nans_halted}
\leanok
\uses{Complexity.LogProg.NAns}
\difficulty{3}
For any nondeterministic register machine, oracle, input, register valuation $v$,
recorded answer $r$, natural number $N$, and Boolean value $b$, a length-$N$ choice
sequence from the halted configuration $(\mathrm{none},v,r)$ leaves recorded answer $b$
if and only if $r=\mathrm{some}\ b$.
\end{lemma}

\begin{lemma}[Answer After a Deterministic Instruction]
\lean{Complexity.LogProg.nans_det}
\label{Complexity.LogProg.nans_det}
\leanok
\uses{Complexity.LogProg.Ins, Complexity.LogProg.NARM.detPart, Complexity.LogProg.NAns, Complexity.LogProg.astep, Complexity.LogProg.narun_cons, Complexity.LogProg.nastep_det}
\difficulty{4}
Let $A$ be a nondeterministic register machine whose instruction at state $l$ is the
deterministic instruction $i$. For any Boolean value $\beta$, register valuation $v$,
recorded answer $r$, natural number $N$, and Boolean value $b$, a length-$(N+1)$ choice
sequence that leaves answer $b$ from $(l,v,r)$ exists if and only if a length-$N$ choice
sequence that leaves answer $b$ exists from the configuration obtained by one
deterministic step using the machine that follows $A$'s deterministic instructions and
replaces each choice instruction by an immediate return of $\beta$.
\end{lemma}

\begin{lemma}[Answer After a Choice Instruction]
\lean{Complexity.LogProg.nans_choose}
\label{Complexity.LogProg.nans_choose}
\leanok
\uses{Complexity.LogProg.NAns, Complexity.LogProg.narun_cons, Complexity.LogProg.nastep_choose}
\difficulty{4}
Let $A$ be a nondeterministic register machine, and suppose its instruction at state $l$
chooses between states $l_1$ and $l_0$. For any register valuation $v$, recorded answer
$r$, natural number $N$, and Boolean value $b$, there is a length-$N+1$ choice sequence
that leaves answer $b$ from the configuration $(l,v,r)$ if and only if such a length-$N$
sequence exists from at least one of $(l_1,v,r)$ or $(l_0,v,r)$.
\end{lemma}

\begin{lemma}[Answer Across a Deterministic Run]
\lean{Complexity.LogProg.nans_run}
\label{Complexity.LogProg.nans_run}
\leanok
\uses{Complexity.LogProg.AConf, Complexity.LogProg.NARM.detPart, Complexity.LogProg.NAns, Complexity.LogProg.NIns.targets, Complexity.LogProg.arun, Complexity.LogProg.narun, Complexity.LogProg.narun_append, Complexity.LogProg.narun_eq_arun}
\difficulty{4}
Let $A$ be a nondeterministic register machine, let $\beta$ be a Boolean value, and
suppose the deterministic machine that follows $A$'s deterministic instructions and
replaces each choice instruction by an immediate return of $\beta$ takes $T$ steps from
configuration $a$ to configuration $a'$. Assume that at every time $t<T$, whenever this
run is still active, the instruction of $A$ at its current state is deterministic. For
any natural number $N$ and Boolean value $b$, a choice sequence of length $T+N$ leaving
answer $b$ from $a$ exists if and only if a choice sequence of length $N$ leaving answer
$b$ from $a'$ exists.
\end{lemma}

\begin{lemma}[Monotonicity of Answer Length]
\lean{Complexity.LogProg.nans_mono_len}
\label{Complexity.LogProg.nans_mono_len}
\leanok
\uses{Complexity.LogProg.AConf, Complexity.LogProg.NAns, Complexity.LogProg.NHalts, Complexity.LogProg.narun_append_of_halted}
\difficulty{3}
Let $A$ be a nondeterministic register machine, with oracle, input $x$, state predicate
$G$, initial configuration $a$, and natural numbers $N\le N'$. Suppose every choice
sequence of length $N$ from $a$ halts, and every earlier prefix configuration that is
still active satisfies $G$. Then a length-$N'$ choice sequence from $a$ that leaves
answer $b$ exists if and only if a length-$N$ choice sequence from $a$ that leaves
answer $b$ exists, for either Boolean value $b$.
\end{lemma}
